\documentclass[12pt]{article}
\usepackage[utf8]{inputenc}
\usepackage[T1]{fontenc}
\usepackage[margin=1in]{geometry}
\usepackage{setspace}
\usepackage{xcolor}
\usepackage[colorlinks=true, urlcolor=blue!60!black, linkcolor=blue!60!black]{hyperref}
\setlength{\parindent}{2em}
\setlength{\parskip}{1.25em}
\renewcommand{\baselinestretch}{1.0}
\begin{document}
\begin{center}
{\LARGE \textit{Marbury v. Madison}, 5 U.S. (1 Cranch) 137 (1803)\par}
\vspace{0.8em}
{\large Case Brief --- Four-Page Edition\par}
\end{center}
\vspace{1em}

\section{Memory Jogger}

Federal courts have the power and duty of judicial review to refuse to enforce acts of Congress that conflict with the Constitution, and Congress cannot expand the Supreme Court's original jurisdiction beyond the categories enumerated in Article III, Section 2 (\href{https://supreme.justia.com/cases/federal/us/5/137/}{Justia}).

\section{Detailed Case Facts}

In the election of 1800, Thomas Jefferson defeated John Adams, and Democratic-Republicans won Congress. In the lame-duck period, the Federalist Congress passed the Judiciary Act of 1801 and the Organic Act of the District of Columbia, which authorized the President to appoint justices of the peace ``to continue in office for five years.'' On March 2--3, 1801, Adams nominated forty-two justices of the peace, including William Marbury. The Senate confirmed them, Adams signed the commissions, and Secretary of State John Marshall, already appointed Chief Justice, affixed the Great Seal. Several commissions, including Marbury's, were still undelivered when Jefferson took office on March 4. Treating them as void, Jefferson told Secretary of State James Madison to withhold them.

\section{Procedural History}

In December 1801, Marbury and three co-petitioners, represented by former Attorney General Charles Lee, filed an original petition in the Supreme Court for a writ of mandamus compelling Madison to deliver the commissions. They relied on Section 13 of the Judiciary Act of 1789, which authorized mandamus ``to any courts appointed, or persons holding office, under the authority of the United States.'' The Court issued a rule to show cause, but Congress repealed the 1801 Act and cancelled the Court's 1802 terms. At argument in February 1803, no one appeared for Madison. The Court decided the case on February 24, 1803.

\section{Judicial Votes}

Unanimous, \textbf{4--0}.
\begin{itemize}
\item \textbf{Majority:} Chief Justice Marshall, joined by Justices Paterson, Chase, and Washington.
\item \textbf{Concurrences and dissents:} None.
\item \textbf{Not participating:} Justices Cushing and Moore (illness).
\end{itemize}

\section{Holding}

Rule discharged; petition dismissed for lack of original jurisdiction. (1) An act of Congress repugnant to the Constitution is void, and federal courts have the authority and duty to refuse to apply it in a case before them. (2) Article III, Section 2 sets an exclusive ceiling on the Court's original jurisdiction, so Section 13 was unconstitutional insofar as it granted original jurisdiction to issue mandamus to executive officers. Marshall's rulings that Marbury had a vested right and a remedy are classically viewed as dicta, because a court without jurisdiction ordinarily cannot reach the merits.

\section{Analysis of Opinions}

\textbf{Majority (Marshall): the right.} Inverting the usual order, Marshall first asked whether Marbury had a right to the commission. Once the President signed it, the appointment was complete, and the seal was conclusive evidence; delivery was a convenience, not a prerequisite. Because the term was fixed at five years, Marbury's right had vested.

\textbf{The remedy.} Where there is a legal right there is a legal remedy. Marshall distinguished discretionary political acts, which are ``only politically examinable,'' from ministerial duties on which individual rights depend, for which officers are ``amenable to the laws.''

\textbf{Jurisdiction and judicial review.} Mandamus was the proper writ, but Article III gives original jurisdiction only in cases ``affecting ambassadors, other public ministers and consuls, and those in which a state shall be a party.'' Letting Congress add to that list would make the distribution ``mere surplusage.'' Facing a conflict between Section 13 and the Constitution, Marshall reasoned that a written constitution is ``the fundamental and paramount law of the nation''; that courts must say what the law is; that Article III extends to cases ``arising under the Constitution''; that the judicial oath would be ``worse than solemn mockery'' if judges enforced unconstitutional laws; and that the Supremacy Clause covers only laws made ``in pursuance of the Constitution.'' Section 13 was void.

\textbf{Separate opinions.} None.

\section{Examples: Future Applications}

\begin{enumerate}
\item \textbf{Same-Side: Enlarged Original Jurisdiction.} Congress gives the Supreme Court original jurisdiction over voters' election-law suits against the Attorney General. No State or diplomat is a party, so the grant is void.
\item \textbf{Same-Side: Treason on One Witness.} A statute permits treason convictions on one witness's testimony, Marshall's own example. It conflicts with Article III's two-witness rule, and courts must vacate the conviction.
\item \textbf{Opposite-Side: District Court Mandamus.} Congress vests district courts with mandamus jurisdiction over cabinet officers' nondiscretionary duties. That does not enlarge the Supreme Court's original jurisdiction.
\item \textbf{Opposite-Side: Recognizing a Foreign Republic.} A citizen seeks mandamus ordering recognition of a breakaway republic. Recognition is a discretionary political power, so courts will refuse.
\item \textbf{Fence-Sitter: Departmentalism.} After the Court strikes a statute, the President complies in that case but enforces an analogous provision against non-parties. \textit{Marbury} establishes judicial review, but scholars debate whether it entails judicial supremacy.
\end{enumerate}

\section{Critique}

\textbf{Doctrinal critique.} William W. Van Alstyne showed that nearly every step was contestable: Marshall arguably should have recused himself, lectured Jefferson on the merits before dismissing for want of jurisdiction, could have read Section 13 narrowly, and offered textual arguments that beg the question of who has final interpretive authority. Originalists such as Prakash, Yoo, and Treanor nonetheless show that judicial review was recognized before 1803.

\textbf{The countermajoritarian difficulty.} Alexander Bickel's phrase captures the tension of unelected judges voiding democratic laws. Sylvia Snowiss and popular constitutionalists argue that modern judicial supremacy far exceeds Marshall's narrower, defensive review; defenders see independent review as indispensable to minority rights and the rule of law.

\section{Key Quotations}

\begin{quote}
``It is emphatically the province and duty of the Judicial Department to say what the law is. Those who apply the rule to particular cases must, of necessity, expound and interpret that rule. If two laws conflict with each other, the Courts must decide on the operation of each.''
\end{quote}
\noindent --- Chief Justice Marshall, for the Court, 5 U.S. (1 Cranch) at 177.

\begin{quote}
``The very essence of civil liberty certainly consists in the right of every individual to claim the protection of the laws whenever he receives an injury.''
\end{quote}
\noindent --- Chief Justice Marshall, for the Court, 5 U.S. (1 Cranch) at 163.

\begin{quote}
``The province of the Court is solely to decide on the rights of individuals, not to enquire how the Executive or Executive officers perform duties in which they have a discretion. Questions, in their nature political or which are, by the Constitution and laws, submitted to the Executive, can never be made in this court.''
\end{quote}
\noindent --- Chief Justice Marshall, for the Court, 5 U.S. (1 Cranch) at 170.

\end{document}
